\documentclass[11pt]{article}
\usepackage[utf8]{inputenc}

%----- Configuración del estilo del documento------%
\usepackage{epsfig,graphicx}
\usepackage[left=2cm,right=2cm,top=1.8cm,bottom=2.3cm]{geometry}
\usepackage{fancyhdr}
\usepackage{lastpage}
\usepackage{url}
\usepackage{longtable}
\pagestyle{fancy}
\fancyhf{}
\fancyhead[R]{\ifnum\value{page}>1 \textbf{\'Algebra L\'ineal II}\fi}
\rfoot{\textit{Página \thepage}}

% Encabezado izquierdo: logo
%\lfoot{\includegraphics[height=1.0cm]{Logo_FC.png}}



%------ Paquetes matemáticos básicos --------%
\usepackage{amsmath}
\usepackage{amssymb}
\usepackage{amsthm}

\usepackage[spanish]{babel}
\usepackage{graphicx}
\usepackage{hyperref}

\usepackage{tabularx}
\usepackage{xcolor}
\usepackage[table]{xcolor}
\usepackage{colortbl}
\usepackage{array, multirow, multicol, tabularx}
\usepackage{tcolorbox}
\newtheorem{theorem}{Theorem}[section]
\newtheorem{corollary}{Corollary}[theorem]
\newtheorem{lemma}[theorem]{Lemma}

%------si-------%
\definecolor{B}{HTML}{FFFFFF}
\definecolor{G}{HTML}{5e5e5e}
\definecolor{R2}{HTML}{d53d40}
\definecolor{A2}{HTML}{034190}
\definecolor{V2}{HTML}{7faa50}

\newcommand{\PrdI}[2]{\ensuremath{\langle #1, #2 \rangle}}
\newcommand{\R}{\mathbb{R}}
\newcommand{\C}{\mathbb{C}}
\newcommand{\N}{\mathbb{N}}
\newcommand{\Z}{\mathbb{Z}}
\newcommand{\Q}{\mathbb{Q}}
\renewcommand{\theenumi}{\Roman{enumi}}
\renewcommand{\labelenumi}{{\theenumi}.}

\begin{document}
	\makeatletter
        \renewenvironment{proof}[1][\proofname]{\par
            \pushQED{\qed}%
            \normalfont \topsep6\p@\@plus6\p@\relax
            \trivlist
            \item\relax
            {\itshape
            #1\@addpunct{.}}\par\vspace{\baselineskip}\ignorespaces
            }{%
            \popQED\endtrivlist\@endpefalse
            }
        \newenvironment{Solucion}[1][Soluci\'on]{\par
            \pushQED{\qed}%
            \normalfont \topsep6\p@\@plus6\p@\relax
            \trivlist
            \item\relax
            {\itshape
            #1\@addpunct{.}}\par\vspace{\baselineskip}\ignorespaces
            }{%
            \popQED\endtrivlist\@endpefalse
            }
            
    \makeatother

%------ Encabezado -------- %

\begin{center}
    \begin{minipage}{3cm}
    	\begin{center}
    		%\includegraphics[height=3.4cm]{logo_unam.png}
    	\end{center}
    \end{minipage}\hfill
    \begin{minipage}{10cm}
    	\begin{center}
    	\textbf{\large Universidad Nacional Autónoma de México}\\[0.1cm]
        \textbf{Facultad de Ciencias}\\[0.1cm]
        \textbf{\'Algebra L\'ineal II}\\[0.1cm]
         Tarea examen 3 \\[0.1cm]
         El\'ias L\'opez Rivera$^{1}$\,\,Adolfo Angel Cardoso Vasquez$^{2}$\\[0.1cm]
        \texttt{\{$^{1}$ elias.lopezr,\,${^2}$ acardosov2400\},@ciencias.unam.mx }\\[0.1cm]
        Fecha:\,\,27/11/2024
    	\end{center}
    \end{minipage}\hfill
    \begin{minipage}{3cm}
    	\begin{center}
    		%\includegraphics[height=3.4cm]{Logo_FC.png}
    	\end{center}
    \end{minipage}
\end{center}

\rule{17cm}{0.1mm}

%------ Fin de encabezado -------- %
\begin{tcolorbox}[title = \textcolor{black}{\textcolor{white}{Problema 1}}]
\textit{Sea $T:V\rightarrow V$ un endomorfismo en un espacio vectorial $V$ de dimensi\'on $2$ sobre el campo $\mathbb{F}$, demuestra que si $T$ no es diagonalizable,entonces $V$ es $T$-ciclico
}
\end{tcolorbox}
\begin{proof}
  Sabemos que $T$ no es diagonalizable entonces no existe una base de vectores propios de $T$ para $V$, por tanto sea $\{v_1,v_2\}$ base de $V$ tenemos que al menos uno de los vectores no es vector propio de $T$, sin perdida de generalidad
  lo es $v_1$, como $v_1\neq 0$, notemos que $T(v_1)\neq 0$, pues de otra manera $v_1$ seria vector propio para el valor propio $0$, luego se tiene que $\{v_1,T(v_1)\}$ es l.i, pues de otra manera $v_1$ seria vector propio de $T$, es decir $\{v_1,T(v_1)\}$ es base de $V$, $V$ es $T$-cilcico
\end{proof}
\begin{tcolorbox}[title = \textcolor{black}{\textcolor{white}{Problema 2}}]
\textit{
Sean \( V \) un espacio vectorial y \( T : V \rightarrow V \) un operador lineal. Sea \( W \) un subespacio \( T \)-invariante de \( V \), y sea \( K \) un subespacio de \( W \) que es \( T|_W \)-invariante (es decir, \( T(W) \subseteq W \) y \( T|_W(K) \subseteq K \)). Demuestra que \( K \) es \( T \)-invariante en \( V \).
}
\end{tcolorbox}
\begin{proof}
    Como $W$ es $T$-invariante $T_{W}(x)=T(x)$ para todo $x\in W$, en particular como $K\subset W$, se sigue que $T(K)=T_{W}(K)$, por tanto $T_{W}(K)=T(K)\subset K$, por tanto $K$ es T-invariante
\end{proof}
\newpage
\begin{tcolorbox}[title = \textcolor{black}{\textcolor{white}{Problema 3}}]
\textit{
Sea \( T : \mathbb{R}^3 \rightarrow \mathbb{R}^3 \) definido por
\[
T(x, y, z) = (y, z, 0).
\]
\begin{itemize}
  \item[(a)] Calcula \( T(v), T^2(v), T^3(v) \) para \( v = (1, 0, 0) \).
  \item[(b)] Determina el subespacio generado por \( \{v, T(v), T^2(v)\} \).
  \item[(c)] Concluye si dicho subespacio es \( T \)-cíclico.
\end{itemize}
}
\end{tcolorbox}
\begin{Solucion}
Notemos que \[ \forall (x,y,z) \in \R^3\] se tiene que 
\begin{align*}
  T^2(x,y,z)=T(T(x,y,z))=T(y,z,0)=(z,0,0)\\
  T^3(x,y,z)=T(T^2(x,y,z))=T(z,0,0)=(0,0,0).
\end{align*}
En especial para \( v = (1, 0, 0) \) se tiene que \(T^k(v)=\vec{0}\; \forall k\geq 1 \ \), por lo que el subespacio \(W\) generado por \( \{v, T(v), T^2(v)\} \) es
\[\mathfrak{L} \{ v, T(v),T^2(v),... \}=\mathfrak{L}_T\{v\}  \]
Luego \(W\) es \( T \)-cíclico.
\end{Solucion}
\begin{tcolorbox}[title = \textcolor{black}{\textcolor{white}{Problema 4}}]
\textit{
Sea
\[
A = \begin{pmatrix} 2 & 1 & 0 \\ 0 & 2 & 1 \\ 0 & 0 & 2 \end{pmatrix}.
\]
\begin{itemize}
  \item[(a)] Encuentra el polinomio característico de \( A \).
  \item[(b)] Calcula los espacios propios \( E_2 \) y \( K_2 = N((A - 2I)^3) \).
  \item[(c)] Determina si \( A \) es diagonalizable.
\end{itemize}
}
\end{tcolorbox}
\begin{Solucion}
El polinomio característico de \( A \) es
\begin{align*}
  \chi_A(t)&=\det(A-tI_3)\\&=\det\begin{pmatrix} 2-t & 1 & 0 \\ 0 & 2-t & 1 \\ 0 & 0 & 2-t \end{pmatrix} \\
  &=(2-t)\det\begin{pmatrix} 2-t & 1 \\ 0 & 2-t \end{pmatrix}\\ &=(2-t)^3. \\
\end{align*}
De aqui se sigue que el polinomio mínimo es de la forma \( \mu_A(t)=(t-2)^k \) con \( 1\leq k \leq 3 \).
Ahora, tenemos que
\[
A-2I_3=\begin{pmatrix}0&1&0\\0&0&1\\0&0&0\\ \end{pmatrix}
\]
\[
(A-2I_3)^2=\begin{pmatrix}0&0&1\\0&0&0\\0&0&0\\ \end{pmatrix}
\]
\[
(A-2I_3)^3=\begin{pmatrix}0&0&0\\0&0&0\\0&0&0\\ \end{pmatrix}
\]
Entonces el polinomio mínimo es \( \mu_A(t)=(t-2)^3 \), y por lo tanto \( A \) no es diagonalizable.\\
Calculamos los espacios propios: \\
El espacio propio \( E_2 \) es
\[E_2=N(A-2I_3)=\left\lbrace \begin{pmatrix}x\\y\\z \end{pmatrix}\in \R^3 : (A-2I_3)\begin{pmatrix}x\\y\\z \end{pmatrix}=
\begin{pmatrix}y\\z\\0 \end{pmatrix}=\begin{pmatrix}0\\0\\0 \end{pmatrix}\right\rbrace 
=\left\lbrace \begin{pmatrix}x\\0\\0 \end{pmatrix}: x\in \R \right\rbrace 
=\mathfrak{L}\left\lbrace \begin{pmatrix}1\\0\\0 \end{pmatrix} \right\rbrace  \]\\
El espacio \( K_2 \) es
\[K_2=N((A-2I_3)^3)=\left\lbrace \begin{pmatrix} x\\y\\z \end{pmatrix}\in \R^3 : \begin{pmatrix}0&0&0\\0&0&0\\0&0&0\\ \end{pmatrix}\begin{pmatrix}x\\y\\z \end{pmatrix}=\begin{pmatrix}0\\0\\0 \end{pmatrix}\right\rbrace =\R^3 \]
\end{Solucion}
\begin{tcolorbox}[title = \textcolor{black}{\textcolor{white}{Problema 5}}]
\textit{
Sea \( T : \mathbb{R}^4 \rightarrow \mathbb{R}^4 \) definido por
\[
[T]_{\beta} = \begin{pmatrix}
0 & 1 & 0 & 0 \\
0 & 0 & 1 & 0 \\
0 & 0 & 0 & 1 \\
0 & 0 & 0 & 0
\end{pmatrix}.
\]
\begin{itemize}
  \item[(a)] Calcula \( T^k(e_1) \) para \( k = 0, 1, 2, 3, 4 \).
  \item[(b)] Determina \( W = \{e_1, T(e_1), T^2(e_1), T^3(e_1)\} \).
  \item[(c)] Verifica que \( W \) es \( T \)-invariante y \( T \)-cíclico, y encuentra el polinomio mínimo de \( T \).
\end{itemize}
}
\end{tcolorbox}
\begin{Solucion}
  Tenemos que \[T(\hat{e}_1)=\vec{0},\] se sigue que \[T^k(\hat{e}_1)=\vec{0}\ \forall k\geq 1,\] por lo que
  \[W=\mathfrak{L}\{\hat{e}_1,T(\hat{e}_1),T^2(\hat{e}_1),T^3(\hat{e}_1)\}=
  \mathfrak{L} \{\hat{e}_1,T(\hat{e}_1),T^2(\hat{e}_1),T^3(\hat{e}_1),\dots\}=\mathfrak{L}_T{ \{\hat{e}_1\}}. \]
  Entonces \( W \) es \( T \)-cíclico y por tanto \(T-\) invariante, esto ultimo se puede verficar del hecho que 
  \( T^k(\hat{e}_1)=\vec{0} \in W\ \forall k\geq 1\), así \(w\in W \Rightarrow w=a\hat{e}_1, \; a\in \R\), de modo que
  \(T(w)=aT(\hat{e}_1)=\vec{0}\in W\).\\
  Finalmente, para calcular el polinomio mínimo de \( T \), notemos que
  \[\chi_T(t)=\det([T]_{\beta}-tI_4)=\det\begin{pmatrix} -t & 1 & 0 & 0 \\ 0 & -t & 1 & 0 \\ 0 & 0 & -t & 1 \\ 0 & 0 & 0 & -t \end{pmatrix}=(-t)^4=t^4.\]
  Por lo tanto el polinomio mínimo es \( \mu_T(t)=t^k,\quad 1\leq k \leq 4 \). Pero \(T(x,y,z,w)=(y,z,w,0)\), asi que \(T^2(x,y,z,w)=(z,w,0,0),\quad T^3(x,y,z,w)=(w,0,0,0)\) y
  \(T^4(x,y,z,w)=(0,0,0,0)\). Entonces el polinomio mínimo es \( \mu_T(t)=t^4 \).

\end{Solucion}
\end{document}